\chapter*{Аннотация}
В работе исследуется, как распределение спектральной энергии по пространственным частотам меняется с глубиной в трёх семействах моделей компьютерного зрения: сверточной ResNet, визуальном трансформере ViT и иерархическом трансформере Swin.

Центральный вклад работы — разработка единого протокола спектрального анализа, впервые позволяющего корректно и непротиворечиво сравнивать разнородные архитектуры. В отличие от предыдущих исследований, протокол явно различает радиальную среднюю мощность (для визуализации) и суммарную энергию колец (для метрик), а также включает контрольные эксперименты, отделяющие эволюцию признаков от тривиального эффекта уменьшения разрешения.

Эксперименты на ImageNet 10K показали: ResNet демонстрирует наиболее устойчивое низкочастотное сжатие; ViT на фиксированной решётке патчей ведёт себя немонотонно; Swin приходит к низкочастотному финальному профилю, но во многом за счёт иерархического понижения разрешения.

Научная новизна состоит в том, что впервые проведено прямое сопоставление спектральной динамики трёх принципиально разных типов архитектур в едином методическом ключе. Предложенный протокол может служить самостоятельным диагностическим инструментом при выборе архитектуры для задач, чувствительных к пространственной структуре признаков.
